\section{EVALUATION AND INFERENCE}\label{ch:eval}

This chapter defines the loss, the estimand, the test, the decision rule and the power study, and states which results are pre-specified and which are not.

\subsection{The loss function}\label{sec:loss}

A forecast of $(\VaR_t,\ES_t)$ is scored against the realised return $y_t$ with the loss
\begin{equation}
\FZ(y_t,\VaR_t,\ES_t;\alpha)=-\frac{\mathbf{1}\{y_t\le \VaR_t\}\,(\VaR_t-y_t)}{\alpha\,\ES_t}+\frac{\VaR_t}{\ES_t}+\log(-\ES_t)-1,
\end{equation}
where $\VaR_t<0$ and $\ES_t<0$ are return quantiles. Lower is better. This is the member of the Fissler--Ziegel family \citep{fissler2016} that \citet{patton2019} call FZ0. Joint elicitability means that the expected loss is smallest at the true VaR and ES pair, so a forecaster cannot lower the expected loss by misreporting either; this is the reason to use such a loss. Scaling the return and both forecasts by $c>0$ adds $\log c$ to the loss, so a difference between the losses of two models does not depend on the units of the return. Section~\ref{sec:sesoi} translates a loss difference into a proportional error in the forecasts.

\subsection{Estimand}\label{sec:estimand}

Let model R (richer) and model B (baseline) be two forecasting procedures, each a model together with its window and refit schedule (Section~\ref{sec:walkforward}). Define the loss difference on forecast day $t$ as
\begin{equation}
d_t=\FZ_t^{R}-\FZ_t^{B},
\end{equation}
so that a negative $d_t$ favours the richer model. This sign convention, richer minus baseline with negative meaning better, is used in every table. The estimand is the long-run mean
\begin{equation}
\mu=\E[d_t],
\end{equation}
the expected loss difference between the two procedures, with the estimation error in their parameters included. It is not a statement about whether the extra parameters of the richer model are zero in a population. This is the sense of equal predictive ability in \citet{giacomini2006}, who compare forecasting methods, including their finite estimation windows, so that the test applies whether or not one model nests the other (see also \citealp{diebold2015}). Their unconditional test is the two-sided Diebold--Mariano statistic; the one-sided version used here is my adaptation. With a rolling window of fixed length, as here (1,260 returns), their asymptotic result remains valid; the counterexample of \citet{mccracken2020} concerns a fixed estimation sample.

\subsection{Estimator and test}\label{sec:test}

The estimator of $\mu$ is the sample mean $\bar d=T^{-1}\sum_{t=1}^{T}d_t$ over the $T$ forecast days. Its standard error $\widehat{\mathrm{SE}}$ comes from a Bartlett-kernel heteroskedasticity- and autocorrelation-consistent (HAC) estimator (Appendix~\ref{app:tech}) with the lag $m=\lfloor 4\,(T/100)^{2/9}\rfloor$, the rule of thumb associated with \citet{newey1994}; at $T=1{,}508$ this gives $m=7$. The test statistic is the Diebold--Mariano statistic \citep{diebold1995},
\begin{equation}
t_{\mathrm{DM}}=\bar d/\widehat{\mathrm{SE}},
\end{equation}
compared with the standard normal distribution; the small-sample correction of \citet{harvey1997} would change it by a factor of $0.9997$ and is not applied. The null hypothesis and its alternative are statements about $\mu$:
\begin{equation}
H_0:\ \mu\ge 0\qquad\text{against}\qquad H_1:\ \mu<0.
\end{equation}
In words, under $H_0$ the richer model does not have a lower expected loss than the baseline. The one-sided $p$-value is $\Phi(t_{\mathrm{DM}})$ and the two-sided $p$-value is $2\{1-\Phi(|t_{\mathrm{DM}}|)\}$. Every table gives the two-sided 95\% interval $\bar d\pm1.96\,\widehat{\mathrm{SE}}$ and labels each $p$-value one- or two-sided.

Normal critical values need not be accurate when one model nests the other \citep{clark2001,clark2007}, but those results concern linear models under squared-error loss and do not carry over to the mixture forecasts and loss used here. For a rolling window, \citet{clark2007} themselves read the usual $t$-statistic as a statement about the estimated-parameter mean difference, which is the reading used here. I therefore read the size of the test from simulation (Section~\ref{sec:power-design}) and not from asymptotic theory.

\subsection{Decision rule and interpretation}\label{sec:decision-rule}

The confirmatory test is RSC\_P against CCC, at tail level $\alpha=5\%$ with a Gaussian chain (Section~\ref{sec:ladder-design}). Two levels should not be confused: $\alpha=5\%$ is the tail level of the forecast, and the test has its own size, also 5\%. The 5\% tail level was chosen for power, since a 5\% VaR is breached about twice as often as a 2.5\% one and gives the loss difference more information; the 2.5\% level of the Basel expected-shortfall standard is the secondary level. The decision rule is one-sided: reject $H_0$ if $t_{\mathrm{DM}}<-1.6449$. I chose a one-sided test because the substantive question has one direction, and because the simulated power and size of Chapter~\ref{ch:results} were computed for exactly this rule. A one-sided test says nothing about harm, so I declare a separate statement: the switching model is called significantly worse if $t_{\mathrm{DM}}>+1.96$. The threshold is 1.96 rather than 1.6449 because in the simulated null the upper-tail rate is 5.2\% at 1.96 but 8.0\% at 1.6449 (Table~\ref{tab:power1508}).

The four possible outcomes are checked in this order and each licenses one sentence, fixed before the run. The full wording of O1, O3 and O4 is in Appendix~\ref{app:wording}.
\begin{description}
\item[O1] $t_{\mathrm{DM}}\le-1.6449$: statistical evidence at the one-sided 5\% level that switching lowers expected FZ0 loss.
\item[O4] $t_{\mathrm{DM}}\ge+1.96$: evidence that the switching specification has higher expected FZ0 loss.
\item[O3] otherwise, and $\bar d-1.6449\,\widehat{\mathrm{SE}}>-0.002$: the data are inconsistent, at the one-sided 5\% level, with an improvement larger than SESOI $=0.002$ FZ0 units.
\item[O2] otherwise: ``The hold-out data neither establish an improvement nor exclude one. The estimate is $\bar d$ (95\% CI $[a,b]$); an improvement of at least $M$ (the minimum detectable difference for the one-sided 5\% test, 80\% power) would have been detected with 80\% probability, and smaller improvements, including none, are consistent with the data.''
\end{description}

No outcome is worded as an absence of an effect. The interval and the minimum detectable difference are reported whatever the outcome.

\subsection{Minimum detectable difference and power}\label{sec:mde}

The minimum detectable difference (MDE) is the smallest true difference that a test detects with 80\% probability, given its standard error. For the two-sided 5\% test it is $(1.96+0.8416)\,\widehat{\mathrm{SE}}=2.8016\,\widehat{\mathrm{SE}}$, and for the one-sided 5\% test it is $(1.6449+0.8416)\,\widehat{\mathrm{SE}}=2.4865\,\widehat{\mathrm{SE}}$. Every comparison table gives the two-sided value, labelled as such, and the primary test gives both. The MDE depends on the design (the sample length and the variability of the loss difference) and not on the estimate, so it can be reported next to a result. I never report power computed at the observed estimate: that quantity is a function of the $p$-value and adds no information \citep{hoenig2001}. Power is reported only for stated effect sizes, which are the simulated scenarios below, the smallest effect of interest and the MDE.

\subsection{The power study}\label{sec:power-design}

I simulate a stylised process with the structure of the problem and score the same pipeline on it. A two-state Markov chain drives an equity-like, a bond-like and a gold-like asset whose volatilities are higher in the stress state and also follow GARCH(1,1) dynamics, with Student-$t$ innovations; Table~\ref{tab:dgp} in Appendix~\ref{app:tech} gives the parameters. The four scenarios differ only in the stress-state correlations. S0 leaves them at their calm values and is the null for the correlation steps; S1 is a moderate breakdown, in which the equity--bond correlation rises from $-0.30$ to $0.00$; S2 is a strong breakdown, to $+0.40$; and S3 is a flight to quality, to $-0.55$. Each replication simulates 1,260 burn-in days and 1,508 test days (the length of the hold-out), runs the whole pipeline of Chapter~\ref{ch:method} including the HMM on the first asset, and applies the test of Section~\ref{sec:test}. There are 500 replications per scenario, with seeds distinct from an earlier study that used 2,520 test days. Size is the rejection rate in S0, and power is the rejection rate in S1 to S3, each with a binomial 95\% interval.

This is an assumption and not a measurement. The parameters are chosen, not estimated from data, and the shift in stress correlation is largest in the stock--bond pair by construction. The choice of RSC\_P as the confirmatory test rested on an earlier run of this simulation, in which switching only the SPY--IEF correlation had equal size and higher power than switching all correlations. That advantage reflects the design of the scenarios and is not a finding about markets. The decision was made before the hold-out was opened and is not reopened here, whatever the 1,508-day study shows.

\subsection{The smallest effect of interest}\label{sec:sesoi}

A non-significant result does not show that an effect is absent. To give a negative result some content I declare one reference value, the smallest effect of interest (SESOI), and use it as a one-sided bound. To express it in economic terms, suppose a forecaster's VaR and ES are both a factor $c$ times the true values of a known distribution. The expected loss then exceeds that of the true forecast by a penalty roughly proportional to $(c-1)^2$; Table~\ref{tab:sesoi} gives it, computed by numerical integration with no market data (\texttt{scripts/confirm\_04\_sesoi.py}). For a normal at $\alpha=5\%$, understating or overstating both forecasts by 3\% costs about 0.0016 to 0.0018, so I declare SESOI $=0.002$, the average of these two penalties to one significant figure. That a 3\% error in a capital-relevant quantity is economically meaningful is my assumption and not a fact. Under a Student-$t$ distribution with six degrees of freedom the penalty is about 20\% lower and rounds to 0.001; I report this as a caveat and keep the declared value.

\begin{table}[t]
\centering\setstretch{1}\footnotesize\setlength{\tabcolsep}{3.5pt}
\caption{\uc{Expected FZ0 penalty of scaling both VaR and ES by a constant factor}}\label{tab:sesoi}
\begin{tabular}{lrrrr}
\toprule
True distribution & $c=0.95$ & $c=0.97$ & $c=1.03$ & $c=1.05$ \\
\midrule
Normal, $\alpha=5\%$ & 0.0051 & 0.0018 & 0.0016 & 0.0042 \\
Student-$t$(6), $\alpha=5\%$ & 0.0040 & 0.0014 & 0.0013 & 0.0034 \\
Normal, $\alpha=2.5\%$ & 0.0067 & 0.0023 & 0.0020 & 0.0055 \\
\bottomrule
\end{tabular}
\tabnote{Analytic expected loss at scale factor $c$ minus the expected loss at $c=1$, for a unit-variance true distribution; no market data are used. Source: \texttt{output/confirm/sesoi.json}.}
\end{table}


The reference value is used in outcome O3 only. With a standard error of 0.0010 to 0.0015, the range suggested by the development data, the one-sided bound $\bar d-1.6449\,\widehat{\mathrm{SE}}$ lies 0.0021 to 0.0025 below the estimate, so O3 requires an estimate near zero or positive and is unlikely; I stated this in advance. A two-one-sided-tests equivalence procedure with margin 0.002 is out of reach, since the half-width of its 90\% interval alone would exceed the margin, so the thesis does not claim equivalence under any outcome.

\subsection{Calibration tests}\label{sec:calibration}

Next to the loss I report, as description only, whether the forecasts have the right coverage. The hit rate, the share of days with $y_t\le\VaR_t$, should be near $\alpha$. The \citet{kupiec1995} test compares the two, the \citet{christoffersen1998} test asks whether hits cluster in time, and their sum is the conditional-coverage test. For expected shortfall I use
\begin{equation}
Z_2=\frac{1}{T\alpha}\sum_{t=1}^{T}\frac{\mathbf{1}\{y_t\le\VaR_t\}\,y_t}{\ES_t}-1,
\end{equation}
which is zero in expectation if the forecast distribution is correct and positive if risk is underestimated. It is the $Z_2$ of \citet{acerbi2014} with the sign reversed, because they write the expected shortfall as a positive loss, and its expectation does not require independent observations. Its $p$-value is the share of 5,000 samples, each drawn from the days' forecast mixtures, whose $Z_2$ is at or above the observed value.

\subsection{Multiple testing and what is pre-specified}\label{sec:prereg}

The confirmatory family has one member, so no adjustment is needed. Besides it I score six secondary analyses of the same forecasts, none of which adds a model, a sample or a regime signal: the six ladder steps at the hold-out; the calendar year 2022 (interval only); a Giacomini--White conditional test of the primary pair with the instruments $\{1,\text{ lagged 63-day SPY--IEF correlation}\}$ (a Wald statistic with a $\chi^2_2$ reference, read with caution because \citet{zhu2020}, in a preprint, show that the conditional null can fail when forecasts rest on estimated parameters); the same comparison at $\alpha=2.5\%$ and in the Student-$t$ chain; a test of superior predictive ability against GARCH-$t$ (\citealp{hansen2005}; stationary bootstrap, 2,000 draws, seed 0; candidates B, D, RSC, RSC\_P); and the calibration tests above. They are unadjusted and labelled secondary, and none is cited as confirmation. Anything else computed is marked POST HOC and carries no inferential weight. The multiple-testing adjustments of the development period \citep{holm1979} are described with those results (Chapter~\ref{ch:results}).

The specification is registered in \texttt{docs/preregistration\_confirmatory.md} and fixed by an annotated git tag together with the scoring code and the raw-data hash, following the logic of pre-registration \citep{nosek2018} within the limits stated in Section~\ref{sec:disclosure}. The run followed the protocol except for one amendment of the scorer's validity rule, made before any scoring (Appendix~\ref{app:disclosure}).
